\section{Introduction}
\label{sec:intro}

Low-light imaging suffers from insufficient illumination, sensor noise,
motion blur, color distortion, and loss of fine structural details.
These effects degrade visual appearance and reduce the reliability of
downstream computer vision systems operating on the captured image.
Restoration methods are commonly optimized using pixel fidelity and
perceptual objectives, such as L1 loss, LPIPS, or structural similarity.
However, improved image-level metrics do not necessarily produce an image
that is more useful for object detection or recognition.

DarkIR is a compact convolutional restoration framework designed for
low-light enhancement and deblurring~\cite{feijoo2025darkir}. It combines
spatial and frequency-domain processing in an asymmetric encoder-decoder
architecture and achieves competitive restoration performance on LOLBlur
and several low-light restoration benchmarks. DarkIR provides a suitable
baseline for this project because its pretrained weights, moderate model
size, and efficient architecture permit controlled fine-tuning and
downstream evaluation.

The central research question is whether task-aware supervision can guide
low-light restoration toward stronger downstream object-detection utility
while maintaining competitive restoration quality. The preliminary project
proposal investigated this problem by proposing a lightweight task-aware
loss based on intermediate representations from a fixed vision model. It
also proposed adaptive feature modulation and a four-way ablation, but these
components were not implemented during the final experiments. The
preliminary report therefore represents a research plan rather than the
description of the final algorithm.

The midterm study implemented a first controlled task-aware experiment,
denoted V1. In V1, a frozen Faster R-CNN ResNet-50 FPN was used to extract
normalized features from restored and paired-reference images~\cite{ren2015fasterrcnn}.
A cosine-distance feature-consistency loss was added to L1 restoration
supervision, while detector parameters remained fixed. The V1
implementation was evaluated after correcting its detector preprocessing
and COCO evaluation protocol. Matched 50,000-step control and task-aware
runs showed no consistent detection improvement; in particular, the
original low-light inputs remained stronger than the enhanced variants
under the frozen COCO-pretrained detector. These results indicated that
feature similarity alone does not guarantee improved detection utility.

Based on this evidence, the final methodology was redesigned as Balanced
V2. Both control and detection-guided arms use the same restoration
objective:
\[
L_{\mathrm{restoration}}
=
L_1
+
0.10L_{\mathrm{LPIPS}}
+
0.05L_{\mathrm{edge}},
\]
and differ only by direct detection-loss supervision. The detection-guided
objective uses the sum of the Faster R-CNN classification, box-regression,
objectness, and RPN box-regression losses from a frozen ExDark-trained
teacher:
\[
L_{\mathrm{V2}}
=
L_{\mathrm{restoration}}
+
\lambda_{\mathrm{det}}L_{\mathrm{detector}}.
\]

Detection weights of 0.05, 0.10, and 0.20 are evaluated. The 0.10 model is
selected as the primary model because it provides the strongest balance of
mAP@50 and mAP@50:95 among the tested guided configurations.

The 5K Balanced V2 experiments show that detection guidance improves
detection relative to the matched control under both an ExDark-trained
teacher and an independent frozen COCO-pretrained detector. The improvement
comes with a modest restoration-quality trade-off. The original low-light
inputs remain stronger in absolute detection scores in both protocols, and
only one training seed is available. The final study therefore evaluates a
conditional restoration--detection trade-off rather than claiming that
image enhancement universally improves detection.
